\documentclass{article}
\usepackage{iclr2027_conference,times}
\input{math_commands.tex}

\usepackage{amsmath,amssymb,amsthm,bm}
\usepackage{booktabs}
\usepackage{graphicx}
\usepackage{microtype}
\usepackage{multirow}
\usepackage{xcolor}
\usepackage{hyperref}
\usepackage{url}

\newtheorem{proposition}{Proposition}
\newcommand{\method}{H-SVDQuant}
\newcommand{\X}{\bm{X}}
\newcommand{\W}{\bm{W}}
\newcommand{\Hess}{\bm{H}}
\newcommand{\D}{\bm{D}}
\newcommand{\Lo}{\bm{L}_1}
\newcommand{\Lt}{\bm{L}_2}
\newcommand{\Qw}{\mathcal{Q}_{W}}
\newcommand{\Qa}{\mathcal{Q}_{A}}
\newcommand{\tr}{\operatorname{tr}}
\newcommand{\pending}{\textcolor{gray}{pending}}
\newcommand{\draftnote}[1]{\textcolor{red!70!black}{[Draft: #1]}}

\title{H-SVDQuant: Hessian-Aware Smoothing and Low-Rank Correction\\
for 4-Bit Large Language Model Inference}

% Keep anonymous for ICLR submission.
\author{Anonymous authors}

% \iclrfinalcopy
\begin{document}
\maketitle

\begin{abstract}
Joint 4-bit weight--activation quantization is attractive for large language
model inference, but activation outliers and anisotropic input sensitivity make
the usual Euclidean smoothing and low-rank corrections unreliable. We introduce
\method{}, a post-training method that optimizes channel smoothing, a
high-precision low-rank branch, and the quantized residual in the geometry
induced by the calibration Hessian. Eliminating the output-side low-rank factor
in closed form yields a deflated Hessian, turning residual quantization into a
standard Hessian-aware problem in the complement of the protected subspace. A
two-channel objective further accounts for online activation-quantization noise.
On Qwen3-8B at W4A4, rank 4, and group size 128, our current implementation
obtains 67.83 MMLU accuracy and 10.50 WikiText-2 perplexity, while an
orthogonal block-Hadamard transform improves these values to 70.11 and 10.18.
We evaluate accuracy independently from systems performance and report operator,
GPU-execution, and synchronized model latency under a common no-offload
protocol. \draftnote{Replace the preliminary numbers and complete the fused
runtime results before submission.}
\end{abstract}

\section{Introduction}

Post-training quantization reduces the memory traffic and arithmetic cost of
large language model (LLM) inference without retraining. Weight-only W4A16
methods such as GPTQ and AWQ are particularly effective for bandwidth-bound
single-token decoding~\citep{frantar2022gptq,lin2023awq}. Quantizing activations
as well can accelerate matrix--matrix products in prompt processing, but 4-bit
activations expose severe channel outliers. SmoothQuant moves outlier magnitude
from activations into weights through an equivalent diagonal
transformation~\citep{xiao2023smoothquant}; rotation methods instead spread
outliers across channels~\citep{ashkboos2024quarot,liu2024spinquant}.

Low-rank-assisted quantization provides a complementary mechanism: preserve a
small, sensitive component in high precision and quantize only the residual.
SVDQuant demonstrates this idea for diffusion transformers and co-designs a
fused W4A4 engine~\citep{li2024svdquant}; LRC studies low-rank correction for
W4A4 LLMs~\citep{scetbon2024lrc}. However, plain truncated SVD treats every
input direction equally. The output error of a linear layer is instead weighted
by the calibration second moment $\Hess=\X^\top\X$. A low-rank direction with
modest Euclidean energy can be critical when the model is highly sensitive along
that direction, while a visually large direction can be nearly irrelevant.

\method{} formulates smoothing, low-rank protection, and residual quantization
in this Hessian geometry. Its central structural result is that, for a fixed
input-side low-rank subspace, the output-side factor can be removed exactly.
The resulting metric is a Hessian with that subspace deflated. This connects
low-rank correction directly to GPTQ-style residual quantization instead of
treating the two components as independent heuristics.

This paper makes three methodological contributions:
\begin{itemize}
    \item We formulate W4A4 low-rank-assisted PTQ with a joint weight-error and
    activation-noise objective in the calibration Hessian geometry.
    \item We derive exact output-factor elimination and Hessian deflation, and
    use these results in a practical alternating solver for smoothing,
    low-rank factors, and integer residual codes.
    \item We isolate the method from orthogonal rotation transforms and from the
    inference backend through factorial accuracy experiments and three-scope
    systems measurements.
\end{itemize}
The fused runtime is an implementation objective, not yet a claimed algorithmic
contribution; it will be promoted to a contribution only after it beats the same
optimized FP16 baseline end to end.

\section{Background and Related Work}

\paragraph{Hessian-aware weight quantization.}
For a linear layer $Y=XW$, second-order PTQ minimizes a layer reconstruction
objective weighted by $H=X^\top X$. GPTQ uses approximate second-order updates
to quantize large language models accurately~\citep{frantar2022gptq}. GANQ
extends this view with GPU-adaptive non-uniform quantization~\citep{zhao2025ganq}.
Our residual-code update uses the same family of solvers, but under a metric
modified by the low-rank branch and activation-noise term.

\paragraph{Smoothing and rotations.}
SmoothQuant transfers activation difficulty into weights using diagonal
rescaling~\citep{xiao2023smoothquant}. QuaRot and SpinQuant use orthogonal
rotations to suppress activation outliers~\citep{ashkboos2024quarot,liu2024spinquant}.
We treat block Hadamard as an external, orthogonal component: ``\method{} +
Hadamard'' is not renamed as our standalone method, and its gain is measured
against the corresponding non-Hessian low-rank baseline with the same transform.

\paragraph{Low-rank-assisted quantization.}
SVDQuant absorbs outliers into a high-precision low-rank branch and fuses that
branch with W4A4 diffusion kernels~\citep{li2024svdquant}. LRC jointly optimizes
quantized weights and low-rank matrices acting on unquantized LLM
activations~\citep{scetbon2024lrc}. Our distinction is geometric: both smoothing
and correction are selected with second-order sensitivity, and low-rank
elimination explicitly changes the metric used to quantize the residual.

\paragraph{Inference systems.}
Marlin is a production-oriented W4A16 kernel for autoregressive
inference~\citep{frantar2024marlin}. QServe shows that dequantization and
non-Tensor-Core work can erase nominal INT4 gains and therefore co-designs the
quantizer and serving kernels~\citep{lin2024qserve}. We follow the same systems
lesson: compression ratio alone is not reported as speedup.

\section{Method}
\label{sec:method}

\subsection{Joint W4A4 formulation}

Consider calibration activations $\X\in\mathbb{R}^{m\times c}$ and a weight
$\W\in\mathbb{R}^{c\times n}$. A positive diagonal smoother $\D$ defines
\begin{equation}
    \widetilde{\X}=\X\D^{-1},\qquad
    \widetilde{\W}=\D\W,
\end{equation}
so the full-precision product is unchanged. We decompose
$\widetilde{\W}=\Lo\Lt+\bm{R}$ and evaluate
\begin{equation}
\widehat{\bm{Y}}=
    \widetilde{\X}\Lo\Lt+
    \Qa(\widetilde{\X})\Qw(\bm{R}).
\label{eq:forward}
\end{equation}
The low-rank factors are stored in FP16/BF16; the residual weights and dynamic
activations use signed 4-bit group quantization.

Let $\bm{E}_A=\Qa(\widetilde{\X})-\widetilde{\X}$ and
$\Delta\W=\bm{R}-\Qw(\bm{R})$. Ignoring the second-order product
$\bm{E}_A\Delta\W$ gives the surrogate
\begin{equation}
\mathcal{L}\approx
\underbrace{\tr(\Delta\W^\top\widetilde{\Hess}\Delta\W)}_{F_W}
+\lambda
\underbrace{\tr(\Qw(\bm{R})^\top\bm{\Sigma}_A\Qw(\bm{R}))}_{F_A},
\label{eq:two-channel}
\end{equation}
where $\widetilde{\Hess}=\D^{-1}\Hess\D^{-1}$ and
$\bm{\Sigma}_A=\mathbb{E}[\bm{E}_A^\top\bm{E}_A]$. $F_W$ measures sensitivity
to residual weight error; $F_A$ penalizes quantized residual energy in channels
where A4 noise is large.

\subsection{Low-rank correction as Hessian deflation}

For fixed $\D$, $\Lo$, and a quantized residual column
$\widehat{\bm w}_j$, the output factor is a small weighted least-squares
problem.

\begin{proposition}[Output-factor elimination]
Let $\bm{M}=\Lo^\top\widetilde{\Hess}\Lo\succ0$. The minimizer of
\[
\min_{\bm\ell\in\mathbb{R}^r}
\|\widetilde{\X}(\widetilde{\bm w}_j-\Lo\bm\ell-
\widehat{\bm w}_j)\|_2^2
\]
is
\begin{equation}
\bm\ell_j^\star=\bm{M}^{-1}\Lo^\top\widetilde{\Hess}
(\widetilde{\bm w}_j-\widehat{\bm w}_j).
\end{equation}
Substitution yields the residual metric
\begin{equation}
\widetilde{\Hess}_{\perp}=
\widetilde{\Hess}-\widetilde{\Hess}\Lo
(\Lo^\top\widetilde{\Hess}\Lo)^{-1}
\Lo^\top\widetilde{\Hess}.
\label{eq:deflation}
\end{equation}
\end{proposition}

Equation~\ref{eq:deflation} has a direct interpretation: the FP branch removes
the directions it spans from the loss seen by the 4-bit residual. Residual
codes can therefore be optimized with a GPTQ/GANQ-style solver after replacing
$\widetilde{\Hess}$ by $\widetilde{\Hess}_{\perp}$.

\subsection{Hessian-aware smoothing and code updates}

We parameterize $\Lo=\D\bm{A}$ so the low-rank branch is evaluated in the
original input coordinates and does not require an additional runtime smoother.
For fixed branch $\bm{B}=\bm{A}\Lt$, we optimize the positive diagonal $\D$
against groupwise models of $F_W+\lambda F_A$. The implementation starts from
the curvature--residual balance
\begin{equation}
d_i\propto
\left(\frac{[\Hess_\perp]_{ii}}
{\|[\W-\bm{B}]_{i,:}\|_p^p+\epsilon}\right)^{1/(p+2)}
\label{eq:d-init}
\end{equation}
and refines $\log d_i$ on cached calibration tokens. Geometric-mean
normalization removes the scale ambiguity.

For fixed $\D$ and $\Lo$, completing the square in
Equation~\ref{eq:two-channel} produces
\begin{equation}
\bm{M}_Q=\widetilde{\Hess}_{\perp}+\lambda\bm{\Sigma}_A,
\qquad
\bm{T}_Q=\bm{M}_Q^{-1}\widetilde{\Hess}_{\perp}\widetilde{\W}.
\label{eq:joint-target}
\end{equation}
We quantize $\bm{T}_Q$ under $\bm{M}_Q$, refit $\Lt$ exactly, and alternate the
blocks. Algorithmic details and derivations are provided in the appendix.

\section{Experimental Design}
\label{sec:experiments}

\paragraph{Separation of claims.}
Accuracy experiments use the eager reference implementation, so they can run
before the fused kernel is complete. Systems experiments use numerically checked
packed backends and compare against the same optimized FP16 implementation.
This prevents a backend limitation from obscuring the quantization method and
prevents an eager implementation from being presented as deployable speedup.

\paragraph{Primary protocol.}
The primary setting is W4A4, low-rank $r=4$, symmetric weight groups of 128,
and per-token activation groups of 128. Group size 128 is a community protocol,
not a method contribution. We calibrate on 128 C4 sequences of length 512 and
evaluate WikiText-2 and C4 perplexity plus MMLU, GSM8K, ARC-Challenge,
ARC-Easy, HellaSwag, and PIQA. The main models are Qwen3-8B and a second 7--8B
model family. \draftnote{Freeze the second model after the first complete
baseline sweep; Llama-3.1-8B is the preferred choice.}

\paragraph{Accuracy comparisons.}
The required rows are FP16, RTN W4A4, GPTQ/AWQ weight-only references,
SmoothQuant, QuaRot or SpinQuant, LRC, adapted SVDQuant, and \method{}. We add a
factorial low-rank comparison:
\begin{center}
\begin{tabular}{lcc}
\toprule
 & No rotation & Block Hadamard \\
\midrule
Non-Hessian low-rank baseline & required & required \\
\method{} & required & required \\
\bottomrule
\end{tabular}
\end{center}
This reports the Hessian gain both without and with the external transform.

\paragraph{Performance comparisons.}
We report: (i) complete Linear operator latency, including smoothing,
activation quantization, residual multiplication, low-rank correction, and
bias; (ii) accumulated GPU execution time in a cached model decode window; and
(iii) synchronized model wall time for the same window. Model loading,
tokenization, network time, and CPU offload are excluded. Prefill and decode are
reported separately with TTFT, TPOT, tokens/s, peak memory, and CUDA event count.

\section{Preliminary Results}

\begin{table*}[t]
\caption{Preliminary Qwen3-8B W4A4 results at rank 4 and group size 128. Hadamard is an
orthogonal existing transform, not part of \method{}. The matched low-rank baselines are
required before these values can support a method comparison. Accuracy is in percent.}
\label{tab:current-accuracy}
\centering
\small
\setlength{\tabcolsep}{4.5pt}
\resizebox{\textwidth}{!}{%
\begin{tabular}{lrrrrrrrr}
\toprule
Method & WT2 PPL $\downarrow$ & MMLU & GSM8K-f & GSM8K-s & ARC-C & ARC-E & HellaSwag & PIQA \\
\midrule
FP16 & 9.72 & 72.98 & 88.02 & 87.49 & 56.57 & 83.54 & 74.92 & 77.69 \\
\method{} & 10.50 & 67.83 & 81.27 & 81.35 & 52.56 & 79.21 & 70.56 & 75.79 \\
\method{} + Hadamard & 10.18 & 70.11 & 85.67 & 85.82 & 53.84 & 80.81 & 72.87 & 76.55 \\
\midrule
Adapted SVDQuant & \pending & \pending & \pending & \pending & \pending & \pending & \pending & \pending \\
Adapted SVDQuant + Hadamard & \pending & \pending & \pending & \pending & \pending & \pending & \pending & \pending \\
LRC & \pending & \pending & \pending & \pending & \pending & \pending & \pending & \pending \\
QuaRot/SpinQuant & \pending & \pending & \pending & \pending & \pending & \pending & \pending & \pending \\
\bottomrule
\end{tabular}
}
\end{table*}


Table~\ref{tab:current-accuracy} establishes the current accuracy anchor but
does not yet identify the Hessian-only gain because the matched non-Hessian
low-rank rows are missing. Hadamard recovers 2.44 points on average across the
seven downstream metrics, but this is evidence of compatibility, not evidence
that Hadamard is part of \method{}.

\begin{table}[t]
\caption{Current Qwen3-8B runtime status on RTX 4090, prompt 256 and decode 32.
Values are diagnostic and are not the proposed fused backend. Speedup is model wall time
relative to the same FP16 run.}
\label{tab:runtime-status}
\centering
\small
\setlength{\tabcolsep}{5pt}
\begin{tabular}{lrrr}
\toprule
Backend & Prefill (ms) & Decode (tok/s) & Speedup \\
\midrule
FP16 dense & 41 & 26.0 & 1.00$\times$ \\
Nunchaku W4A4 & 63 & 16.0 & 0.63$\times$ \\
Hybrid W4A4/W4A16 & 67 & 2.6 & 0.10$\times$ \\
Local W4A16 & 1675 & 2.3 & 0.09$\times$ \\
Local W4A4 WMMA & 418 & 2.5 & 0.09$\times$ \\
\bottomrule
\end{tabular}
\end{table}


The current runtime results in Table~\ref{tab:runtime-status} are diagnostic:
no measured path accelerates the complete model over FP16. Marlin confirms that
production W4A16 can accelerate the large MLP projections, but its roughly
constant launch/tile floor loses on smaller Q/K/V projections. Splitting smooth,
packed multiplication, and rank-4 correction into separate launches then erases
the remaining gain. Our next backend groups QKV and gate/up projections and
fuses smoothing and low-rank correction with the packed epilogue. A systems
claim is accepted only if the fused path passes all three scopes.

\section{Ablations and Analysis}

The main ablation is a $2\times2$ decomposition of the two Hessian mechanisms:
plain smoothing/plain SVD; Hessian smoothing/plain SVD; plain smoothing/Hessian
low rank; and both. We additionally sweep rank $r\in\{0,2,4,8,16\}$, weight and
activation group size, calibration size, alternating iterations, and
$\lambda$. Layerwise $F_W$, $F_A$, output NMSE, and accumulated hidden-state
drift connect the optimization objective to downstream accuracy.

For systems analysis, we sweep matrix rows $M\in\{1,2,4,8,16,64,128,256,1024\}$
for Q/O, K/V, gate/up, and down shapes. We report event counts and a latency tax
decomposition for smoothing, packed multiplication, low-rank down/up
projections, and framework dispatch. This determines whether a failure is an
algorithmic arithmetic tax, a kernel throughput problem, or launch overhead.

\section{Limitations}

The present evidence is limited to one fully evaluated 8B model and one RTX
4090 runtime. The optional Hadamard path is implemented only in the eager
backend, and cached-generation behavior must be validated before reporting its
generation tasks. Low-rank FP factors reduce the nominal compression ratio and
add arithmetic whose relative cost depends on rank, matrix shape, batch size,
and GPU architecture. Finally, the additive-noise surrogate omits a second-order
activation--weight error term; its usefulness is empirical rather than an exact
model of deterministic quantization.

\section{Conclusion}

\method{} places smoothing, low-rank protection, and residual quantization in a
shared Hessian geometry. Exact low-rank factor elimination exposes a deflated
metric that can be consumed by established second-order quantizers, while a
joint activation-noise term targets the failure mode specific to W4A4. The
remaining work is deliberately separated: matched accuracy baselines establish
the algorithmic contribution, and a fused grouped backend must establish any
deployment claim against optimized FP16.

\subsubsection*{AI use statement}
Generative AI tools were used to assist with document organization, language
editing, benchmark-code development, and LaTeX preparation. The authors are
responsible for the method, experiments, verification of all generated code and
claims, and the final text. \draftnote{Coauthors must review this disclosure
against the final ICLR policy and the actual usage before submission.}

\bibliography{references}
\bibliographystyle{iclr2027_conference}

\appendix
\section{Derivation of Hessian deflation}

Writing $\bm e_j=\widetilde{\bm w}_j-\widehat{\bm w}_j$, the fixed-column loss
is
\[
f(\bm\ell)=(\bm e_j-\Lo\bm\ell)^\top
\widetilde{\Hess}(\bm e_j-\Lo\bm\ell).
\]
Setting its gradient to zero gives
$\Lo^\top\widetilde{\Hess}\Lo\bm\ell=
\Lo^\top\widetilde{\Hess}\bm e_j$. Substitution of the solution gives
$\bm e_j^\top\widetilde{\Hess}_\perp\bm e_j$ with
$\widetilde{\Hess}_\perp$ from Equation~\ref{eq:deflation}. Thus the
high-precision branch creates a rank-$r$ null space in the residual
quantization metric.

\section{Alternating calibration algorithm}

For each Linear layer, we accumulate $H=X^\top X$ and a bounded activation
reservoir. Starting from an empty branch, each outer iteration performs:
\begin{enumerate}
    \item deflate $H$ using the previous input-side branch;
    \item initialize and refine $D$ using Equation~\ref{eq:d-init} and the
    groupwise two-channel surrogate;
    \item compute a rank-$r$ weighted approximation under the smoothed Hessian;
    \item form the joint target in Equation~\ref{eq:joint-target} and apply
    GPTQ-style grouped residual quantization;
    \item refit $L_2$ by weighted least squares and retain the candidate with
    the smallest measured calibration output error.
\end{enumerate}
All downstream evaluations use the same saved integer codes, scales, smoother,
and low-rank factors.

\section{Required supplementary tables}

The final supplement will contain per-model task metrics, calibration
sensitivity, rank and group sweeps, cached/no-cache equivalence checks,
per-shape operator latency distributions, Nsight/CUDA event summaries, and
hardware/software versions. Seeds and raw JSON files will be released with the
artifact.

\end{document}
